%! Potential Errors When Performing Tests — Unit 6, Lesson 6.7 — Exit Ticket (Answer Key)
\documentclass[10pt]{article}
\usepackage{apstats-article}
\usepackage{apstats-key}   % pulls in -boxes; do NOT also load -boxes
\newcommand{\TallMath}[1]{$\displaystyle #1\rule[-1.4em]{0pt}{3.2em}$}

\begin{document}

\pageheader{Unit 6, Lesson 6.7}{Exit Ticket --- Answer Key}
\namedateperiod

\vspace{0.15in}

$H_0: p=0.02$; $H_a: p>0.02$.

\begin{enumerate}
  \item \ansline{Type I error: concluding that the proportion of underweight packages exceeds 2\% when it actually does not. The plant would be shut down or processes changed unnecessarily.}
  \item \ansline{Type II error: failing to detect that the proportion of underweight packages truly exceeds 2\%. Underweight packages continue to reach consumers without correction.}
  \item \ansline{Type II error is more serious from a consumer safety perspective --- consumers receive underweight packages (they are shortchanged). Continuing to ship defective products also exposes the company to legal liability.}
  \item \ansline{Power increases --- a larger sample reduces $SE_{\hat p}$, making it easier to detect a true excess defect rate above 2\%.}
\end{enumerate}

\end{document}
